% GENERATED FILE. Edit canonical lessons, then run npm run generate:books.
% canonical-source-hash: fnv1a64:2523e300243f0e76
% canonical-lessons: MR-C537-abhinay, MR-C537-bhumika, MR-C537-talim, MR-C537-vesbhusa, MR-C537-sakarne

\chapter{On Stage}
\label{ch:mr-on-stage}
\hlchaptermodality{537}

\begin{chapteropening}
\textbf{By the end of this chapter:} \emph{I can talk about acting, roles, rehearsals and costumes.}

Complete the last lesson of chapter 537: I can talk about acting, roles, rehearsals and costumes.
\end{chapteropening}

\section[\texorpdfstring{\mr{अभिनय} (abhinay)}{अभिनय (abhinay)}]{\mr{अभिनय} (abhinay) — acting}

\label{lesson:MR-C537-abhinay}



\begin{quote}
Before the new one: say the Marathi for to hum, then the Marathi for to practise music.
\end{quote}

\subsection*{You'll want to know: \mr{अभिनय}}
\textbf{\mr{अभिनय}} — \emph{abhinay} — \textquotedblleft{}acting\textquotedblright{}.

\begin{quote}
\mr{त्या} \mr{अभिनेत्रीचा} \mr{अभिनय} \mr{सगळ्यांना} \mr{आवडला}.
\end{quote}

\begin{quote}
\emph{tyā abhinetrīcā abhinay sagḷyānnā āvaḍlā.}
\end{quote}

\begin{quote}
Everyone liked that actress's acting.
\end{quote}

Masculine.

\begin{cousinweb}
\textbf{\mr{अभिनय}} (\emph{abhinay}) is a tatsama, taken from Sanskrit \emph{abhinaya}, \textquotedblleft{}expressing a feeling on stage, acting\textquotedblright{} (Monier-Williams).
\end{cousinweb}

\begin{grammarlens}[title={What you've built}]
One more word for describing what you lived and felt.
\end{grammarlens}

\subsection*{Your turn}
\begin{itemize}
\raggedright
  \item \emph{Say it:} \emph{abhinay}
  \item \emph{Say it:} \emph{abhinay}, once more
  \item \emph{Say it:} \emph{tyā abhinetrīcā abhinay sagḷyānnā āvaḍlā.} — Everyone liked that actress's acting.
  \item \emph{Recall:} say the Marathi for to hum, then the Marathi for to practise music, then say \emph{abhinay} again
\end{itemize}

\begin{tcolorbox}[breakable,colback=teal!4,colframe=teal!35!black,title={Before you move on}]
What does \mr{अभिनय} mean? (\textquotedblleft{}Acting\textquotedblright{}.) Say it once more.
\end{tcolorbox}
\section[\texorpdfstring{\mr{भूमिका} (bhūmikā)}{भूमिका (bhūmikā)}]{\mr{भूमिका} (bhūmikā) — a role; a part}

\label{lesson:MR-C537-bhumika}



\begin{quote}
Before the new one: say the Marathi for to practise music, then the Marathi for acting.
\end{quote}

\subsection*{You'll want to know: \mr{भूमिका}}
\textbf{\mr{भूमिका}} — \emph{bhūmikā} — \textquotedblleft{}a role; a part\textquotedblright{}.

\begin{quote}
\mr{नाटकात} \mr{अरुणने} \mr{राजाची} \mr{भूमिका} \mr{केली}.
\end{quote}

\begin{quote}
\emph{nāṭkāt aruṇne rājācī bhūmikā kelī.}
\end{quote}

\begin{quote}
Arun played the king in the play.
\end{quote}

Feminine.

\mr{भूमिका} \mr{करणे}, to play a part.

\begin{cousinweb}
\textbf{\mr{भूमिका}} (\emph{bhūmikā}) is a tatsama, taken from Sanskrit \emph{bhūmikā}, \textquotedblleft{}ground\textquotedblright{}, and also \textquotedblleft{}the part an actor takes\textquotedblright{} (Monier-Williams).
\end{cousinweb}

\begin{grammarlens}[title={What you've built}]
One more word for describing what you lived and felt.
\end{grammarlens}

\subsection*{Your turn}
\begin{itemize}
\raggedright
  \item \emph{Say it:} \emph{bhūmikā}
  \item \emph{Say it:} \emph{bhūmikā}, once more
  \item \emph{Say it:} \emph{nāṭkāt aruṇne rājācī bhūmikā kelī.} — Arun played the king in the play.
  \item \emph{Recall:} say the Marathi for to practise music, then the Marathi for acting, then say \emph{bhūmikā} again
\end{itemize}

\begin{tcolorbox}[breakable,colback=teal!4,colframe=teal!35!black,title={Before you move on}]
What does \mr{भूमिका} mean? (\textquotedblleft{}A role; a part\textquotedblright{}.) Say it once more.
\end{tcolorbox}
\section[\texorpdfstring{\mr{तालीम} (tālīm)}{तालीम (tālīm)}]{\mr{तालीम} (tālīm) — a rehearsal}

\label{lesson:MR-C537-talim}



\begin{quote}
Before the new one: say the Marathi for acting, then the Marathi for a role.
\end{quote}

\subsection*{You'll want to know: \mr{तालीम}}
\textbf{\mr{तालीम}} — \emph{tālīm} — \textquotedblleft{}a rehearsal\textquotedblright{}.

\begin{quote}
\mr{नाटकाची} \mr{तालीम} \mr{रोज} \mr{संध्याकाळी} \mr{चालते}.
\end{quote}

\begin{quote}
\emph{nāṭkācī tālīm roj sandhyākāḷī cālte.}
\end{quote}

\begin{quote}
The play's rehearsal goes on every evening.
\end{quote}

Feminine.

\begin{cousinweb}
\textbf{\mr{तालीम}} (\emph{tālīm}) is a loan through Persian usually traced to Arabic \emph{taʿlīm}, \textquotedblleft{}teaching\textquotedblright{}.
\end{cousinweb}

\begin{grammarlens}[title={What you've built}]
One more word for describing what you lived and felt.
\end{grammarlens}

\subsection*{Your turn}
\begin{itemize}
\raggedright
  \item \emph{Say it:} \emph{tālīm}
  \item \emph{Say it:} \emph{tālīm}, once more
  \item \emph{Say it:} \emph{nāṭkācī tālīm roj sandhyākāḷī cālte.} — The play's rehearsal goes on every evening.
  \item \emph{Recall:} say the Marathi for acting, then the Marathi for a role, then say \emph{tālīm} again
\end{itemize}

\begin{tcolorbox}[breakable,colback=teal!4,colframe=teal!35!black,title={Before you move on}]
What does \mr{तालीम} mean? (\textquotedblleft{}A rehearsal\textquotedblright{}.) Say it once more.
\end{tcolorbox}
\section[\texorpdfstring{\mr{वेशभूषा} (veśbhūṣā)}{वेशभूषा (veśbhūṣā)}]{\mr{वेशभूषा} (veśbhūṣā) — a costume}

\label{lesson:MR-C537-vesbhusa}



\begin{quote}
Before the new one: say the Marathi for a role, then the Marathi for a rehearsal.
\end{quote}

\subsection*{You'll want to know: \mr{वेशभूषा}}
\textbf{\mr{वेशभूषा}} — \emph{veśbhūṣā} — \textquotedblleft{}a costume\textquotedblright{}.

\begin{quote}
\mr{नाटकासाठी} \mr{मुलांनी} \mr{जुन्या} \mr{काळची} \mr{वेशभूषा} \mr{केली}.
\end{quote}

\begin{quote}
\emph{nāṭkāsāṭhī mulānnī junyā kāḷcī veśbhūṣā kelī.}
\end{quote}

\begin{quote}
For the play the children dressed in costumes of the old days.
\end{quote}

Feminine.

\mr{वेशभूषा} \mr{करणे}, to put on a costume.

\begin{cousinweb}
\textbf{\mr{वेशभूषा}} (\emph{veśbhūṣā}) is a tatsama compound of Sanskrit \emph{veśa}, \textquotedblleft{}dress\textquotedblright{}, and \emph{bhūṣā}, \textquotedblleft{}ornament\textquotedblright{} (Monier-Williams).
\end{cousinweb}

\begin{grammarlens}[title={What you've built}]
One more word for describing what you lived and felt.
\end{grammarlens}

\subsection*{Your turn}
\begin{itemize}
\raggedright
  \item \emph{Say it:} \emph{veśbhūṣā}
  \item \emph{Say it:} \emph{veśbhūṣā}, once more
  \item \emph{Say it:} \emph{nāṭkāsāṭhī mulānnī junyā kāḷcī veśbhūṣā kelī.} — For the play the children dressed in costumes of the old days.
  \item \emph{Recall:} say the Marathi for a role, then the Marathi for a rehearsal, then say \emph{veśbhūṣā} again
\end{itemize}

\begin{tcolorbox}[breakable,colback=teal!4,colframe=teal!35!black,title={Before you move on}]
What does \mr{वेशभूषा} mean? (\textquotedblleft{}A costume\textquotedblright{}.) Say it once more.
\end{tcolorbox}
\section[\texorpdfstring{\mr{साकारणे} (sākārṇe)}{साकारणे (sākārṇe)}]{\mr{साकारणे} (sākārṇe) — to portray; to bring to life}

\label{lesson:MR-C537-sakarne}



\begin{quote}
Before the new one: say the Marathi for a rehearsal, then the Marathi for a costume.
\end{quote}

\subsection*{You'll want to know: \mr{साकारणे}}
\textbf{\mr{साकारणे}} — \emph{sākārṇe} — \textquotedblleft{}to portray; to bring to life\textquotedblright{}.

\begin{quote}
\mr{तिने} \mr{आईची} \mr{भूमिका} \mr{खूप} \mr{छान} \mr{साकारली}.
\end{quote}

\begin{quote}
\emph{tine āīcī bhūmikā khūp chān sākārlī.}
\end{quote}

\begin{quote}
She portrayed the mother beautifully.
\end{quote}

Past, with -\mr{ने} (he …): \mr{त्याने} \mr{भूमिका} \mr{साकारली}, \emph{tyāne bhūmikā sākārlī}.

\begin{cousinweb}
\textbf{\mr{साकारणे}} (\emph{sākārṇe}) is built on Sanskrit \emph{sākāra}, \textquotedblleft{}having form\textquotedblright{} (Monier-Williams): to give a role a shape.
\end{cousinweb}

\begin{grammarlens}[title={What you've built}]
One more word for describing what you lived and felt.
\end{grammarlens}

\subsection*{Your turn}
\begin{itemize}
\raggedright
  \item \emph{Say it:} \emph{sākārṇe}
  \item \emph{Say it:} \emph{sākārṇe}, once more
  \item \emph{Say it:} \emph{tine āīcī bhūmikā khūp chān sākārlī.} — She portrayed the mother beautifully.
  \item \emph{Recall:} say the Marathi for a rehearsal, then the Marathi for a costume, then say \emph{sākārṇe} again
\end{itemize}

\begin{tcolorbox}[breakable,colback=teal!4,colframe=teal!35!black,title={Before you move on}]
What does \mr{साकारणे} mean? (\textquotedblleft{}To portray; to bring to life\textquotedblright{}.) Say it once more.
\end{tcolorbox}
